% P10-Evaluationsfeinschliff, 16.09.2026: Aussagegrenzen und Anforderungsbilanz; keine neuen Messwerte.
% P10-Core-Inhaltsprüfung, 15.09.2026: datierter Ergebnisanschluss; B-60 korrigiert.
% Belege: thesis-evidence/w2/core-content-audit-20260915/evaluation/
% B-59, 15.09.2026: datierte Ledger-Nachauswertung, Stützung und Reichweite.
% Beleg: thesis-evidence/w2/ledger-repetitions-20260915/ (Originale unverändert).
% M19 (12.09.2026): Leserführung und präzise Ergebnisdarstellung;
% Messwerte und historische Nachweisgrenzen bleiben erhalten.
% Phase 02b (10.09.2026): interpretierende Ledger-Einordnung nach 8.1
% verdichtet; Ergebnisprofil, Tabellen und Aufwandserhebung erhalten.
% ============================================================
% §7.7 Zusammenführung und Aufwandsprofil — M02 (10.09.2026)
% B-14: Input-Stresswert 135.592; Verhältnis Ledger/F rund 1,3.
% NACHREVIEW: Zahlen auf konsolidierten Stand (37/55, 0,495, 53/55);
% neuer Eigenschafts-Absatz (Auftrag §5/7.7: Ledger an eigenen
% belegten Eigenschaften, kein Gesamtsieger); Zusatz-HERKUNFT:
% runs/w2-nachreview-konsolidierung.json.
% NEUE 8er-Struktur (Schreibplan-Zeile 7.7: je ANFORDERUNG Befund ·
% Evidenz · verbleibende Grenze, KEIN Score; Aufwand mit gleicher
% Messgrenze, Tokens inkl. Wiederholungen/Reparaturen; Fehlversuche
% gesondert). Schließt die Prüflandkarte aus 7.1 (P-Nummern
% identisch). Budget ~1,2 S. + 2 Tabellen.
% ============================================================
% HERKUNFT: Befunde/Zahlen = tex-v2 §§7.3–7.6 (dortige Kanonik);
% Token-Zahlen: thesis-evidence/w2/token-usage-correction-20260910.json.
% Letzter kumulativer Wert je Token-Metrikserie, gegengeprüft gegen
% einzelne Chat-Spans; Cache-Anteile und Orchestrierungs-Summen nicht
% zusätzlich addieren. Alle vier Hauptläufe nachgerechnet (B-14).
% Modellreihe (inkl. 4 F-Ausfälle) = Anhang.
% „untersucht, aber nicht erfüllt" vs. „nicht untersucht" bleibt
% unterscheidbar (Umsetzungsplan §7.7).
% ------------------------------------------------------------

\section{Zusammenführung und Aufwandsprofil}
\label{sec:eval-ergebnisprofil}

Tabelle~\ref{t:eval-ergebnisprofil} schließt die Prüflandkarte
aus Abschnitt~\ref{sec:eval-ziele}: Sie verbindet die Prüffragen mit
A1--A10 und führt Befund, tragende Evidenz und verbleibende Grenze
zusammen. Für A1/A2 ergänzt sie die querschnittlichen Ausführungsbelege.
Eine Anforderung kann durch mehrere Prüffragen untersucht werden;
die Zuordnung bedeutet keine pauschale Erfüllung. Bewusst gibt es keinen Gesamtwert
und keine nachträglich definierte Bestehensschwelle --- die
Befunde sind ein mehrdimensionales Profil, und „nicht
untersucht`` bleibt von „untersucht, aber nicht erfüllt``
unterscheidbar.

% P10-Evaluationsfeinschliff, 16.09.2026: Aussagegrenzen und Anforderungsbilanz; keine neuen Messwerte.
% P10-Core-Inhaltsprüfung, 15.09.2026: datierter Ergebnisanschluss; B-60 korrigiert.
% Belege: thesis-evidence/w2/core-content-audit-20260915/evaluation/
% B-59, 15.09.2026: datierte Ledger-Nachauswertung, Stützung und Reichweite.
% Beleg: thesis-evidence/w2/ledger-repetitions-20260915/ (Originale unverändert).
% P10-4a, 14.09.2026: Mechanismus, Quellenreichweite und v4-Nachweis angeglichen.
% M19 (12.09.2026): Leserführung und präzise Ergebnisdarstellung;
% Messwerte und historische Nachweisgrenzen bleiben erhalten.
% ============================================================
% Tabelle 7.7-1: Ergebnisprofil (schließt die Prüflandkarte) — v01
% In Overleaf ablegen als: tables/Chap7-v02/chap7.7table1.tex
% Aufruf im Fließtext: % P10-Evaluationsfeinschliff, 16.09.2026: Aussagegrenzen und Anforderungsbilanz; keine neuen Messwerte.
% P10-Core-Inhaltsprüfung, 15.09.2026: datierter Ergebnisanschluss; B-60 korrigiert.
% Belege: thesis-evidence/w2/core-content-audit-20260915/evaluation/
% B-59, 15.09.2026: datierte Ledger-Nachauswertung, Stützung und Reichweite.
% Beleg: thesis-evidence/w2/ledger-repetitions-20260915/ (Originale unverändert).
% P10-4a, 14.09.2026: Mechanismus, Quellenreichweite und v4-Nachweis angeglichen.
% M19 (12.09.2026): Leserführung und präzise Ergebnisdarstellung;
% Messwerte und historische Nachweisgrenzen bleiben erhalten.
% ============================================================
% Tabelle 7.7-1: Ergebnisprofil (schließt die Prüflandkarte) — v01
% In Overleaf ablegen als: tables/Chap7-v02/chap7.7table1.tex
% Aufruf im Fließtext: % P10-Evaluationsfeinschliff, 16.09.2026: Aussagegrenzen und Anforderungsbilanz; keine neuen Messwerte.
% P10-Core-Inhaltsprüfung, 15.09.2026: datierter Ergebnisanschluss; B-60 korrigiert.
% Belege: thesis-evidence/w2/core-content-audit-20260915/evaluation/
% B-59, 15.09.2026: datierte Ledger-Nachauswertung, Stützung und Reichweite.
% Beleg: thesis-evidence/w2/ledger-repetitions-20260915/ (Originale unverändert).
% P10-4a, 14.09.2026: Mechanismus, Quellenreichweite und v4-Nachweis angeglichen.
% M19 (12.09.2026): Leserführung und präzise Ergebnisdarstellung;
% Messwerte und historische Nachweisgrenzen bleiben erhalten.
% ============================================================
% Tabelle 7.7-1: Ergebnisprofil (schließt die Prüflandkarte) — v01
% In Overleaf ablegen als: tables/Chap7-v02/chap7.7table1.tex
% Aufruf im Fließtext: \input{tables/Chap7-v02/chap7.7table1}
% Label: t:eval-ergebnisprofil
% Stil: chap6.1-Konvention. P-Nummern = t:eval-pruefkarte (7.1).
% HERKUNFT: Befunde = tex-v2 §§7.3–7.6 (Zahlen: Gold v2/Nenner 109;
% Kanonik runs/w2-goldv2-eval-*.json, z1-stufenanalyse.md rev2,
% fehlerklassen-matrix.md v2, z2-fallblaetter.md rev2,
% traceability-audit.json, Lauf dc711f).
% ============================================================

\begin{table}[htbp]
  \centering
  \small
  \begin{tabular}{>{\raggedright\arraybackslash}p{0.08\textwidth}>{\raggedright\arraybackslash}p{0.40\textwidth}>{\raggedright\arraybackslash}p{0.42\textwidth}}
    \hline
    \textbf{Bezug} & \textbf{Befund (Evidenzklasse)} & \textbf{Verbleibende Grenze} \\
    \hline
    A1/A2 & Ausgewählte Verarbeitungsschritte, Übergaben und Zustandswirkungen aus Laufartefakten rekonstruierbar; Nachrichten- und Prozessgrenzentests ergänzen diese Belege [D]+[VA] & Keine Vollprüfung aller Informations- und Protokollierungspfade; sichtbarer Kontext belegt keine Nutzung durch das Modell \\
    \shortstack[l]{P1\\(A3)} & Hauptläufe: Full Coverage 0{,}839 / 0{,}495 (Treue/Stress); Korrektheit jeweils 1{,}0, Stützung 1{,}0 / 0{,}957. Drei Stressfall-Ausgaben: Full Coverage 0{,}495--0{,}523 [B] & Nur 43/109 Aussagen in allen drei Ausgaben full; bekannte Quelle, keine allgemeine Stabilitätsaussage. Im Hauptlauf entstanden zwei der 55 Lücken erst nach der Extraktion; teilweise unvollständiger Referenzkontext \\
    \shortstack[l]{P2\\(A3)} & 37 der 55 Stressfall-Lücken mit zuordenbarem Hinweis (Silent-Miss 0{,}327); Treue-Fall 5/5 still [B] & Prüfpotenzial, keine gemessene Behebung; Hinweise verrauscht (18/48 relevant, 1 explizite Warnung); stille Klasse: Detailverlust in verwendeten Units \\
    \shortstack[l]{P3\\(A4)} & 107/109 Klassen unverändert; 2 Verluste aus einer Zusammenführung; 6/7 Merges ohne Klassenverlust [B] & Übergangs-Check prüft Referenz, nicht Inhalt (belegte Kontrolllücke); Verifikation: 12 Paare, nicht verblindet \\
    \shortstack[l]{P4\\(A4/A6)} & Freigabe-Governance am realen Apply-Pfad getestet (4/4); Challenge-Bilanz 6 getestet / 4 teilweise / 2 offen [VA] & Initialer Seed ohne Einzelautorisierung; G01 (nur Warnung) und G04 (Naht-Umgehung) ungetestet; Prüftiefe überwiegend In-Memory (28/36) \\
    \shortstack[l]{P5\\(A5/A7)} & Fünf Falltypen und acht zusätzliche Inhaltsprüfungen: autorisierte Fortschreibung, passende PBI-/Feature-Verbindungen und Ablehnung einer unvollständigen Vorlage fallbezogen belegt [D] & Historische, nicht unabhängige Fälle; teilweise unvollständige Einordnung und Inhaltserhaltung, begrenzte Folgeangleichung; keine Qualitätsquote \\
    \shortstack[l]{P6\\(A3/A5)} & v4: 104/104 Ledger-Ursprungskanten auflösbar; 234/237 Items mit direktem Bezug oder Beitragspfad; 64 aktuelle Einzelreferenzen bei 33 Items offen [VA] & Teilbeitrag ist keine vollständige Herkunft; getrennte Bezugsgrößen, keine semantische Vollprüfung; ein Snapshot \\
    \shortstack[l]{P7\\(A8)} & Drei Issues real erzeugt; ausgewählter Anforderungsinhalt bis zur Projektion feldverglichen; externer Kommentar-Rückweg bis zum Issue-Update; Wiedererkennung abgelehnter Eingänge [D] & kein allgemeiner Drift-Monitor (G08 teilweise); ein Vermerk-Write zunächst fehlgeschlagen \\
    \shortstack[l]{P8\\(A9/A10)} & Wiederaufnahme ohne erneute Ausführung der abgeschlossenen Ledger-Stufe [D]; separater Prozessgrenzen-Smoke mit 14 Prüfungen [VA]; 2 Steward-Aufgaben über reguläre Gates [D] & keine Fortsetzung \emph{innerhalb} laufender Verarbeitung gezeigt; kein Endzustandsvergleich; keine Usability-Aussage \\
    \hline
  \end{tabular}
  \caption[Ergebnisprofil entlang der Prüffragen]{Ergebnisprofil entlang der Prüflandkarte aus
    Tabelle~\ref{t:eval-pruefkarte} --- Befund, Evidenzklasse und
    verbleibende Grenze je Prüffrage mit Bezug zu A1--A10; A1/A2
    querschnittlich ergänzt. Keine pauschale Erfüllungswertung, kein
    Gesamtwert und keine Bestehensschwelle.}
  \label{t:eval-ergebnisprofil}
\end{table}

% Label: t:eval-ergebnisprofil
% Stil: chap6.1-Konvention. P-Nummern = t:eval-pruefkarte (7.1).
% HERKUNFT: Befunde = tex-v2 §§7.3–7.6 (Zahlen: Gold v2/Nenner 109;
% Kanonik runs/w2-goldv2-eval-*.json, z1-stufenanalyse.md rev2,
% fehlerklassen-matrix.md v2, z2-fallblaetter.md rev2,
% traceability-audit.json, Lauf dc711f).
% ============================================================

\begin{table}[htbp]
  \centering
  \small
  \begin{tabular}{>{\raggedright\arraybackslash}p{0.08\textwidth}>{\raggedright\arraybackslash}p{0.40\textwidth}>{\raggedright\arraybackslash}p{0.42\textwidth}}
    \hline
    \textbf{Bezug} & \textbf{Befund (Evidenzklasse)} & \textbf{Verbleibende Grenze} \\
    \hline
    A1/A2 & Ausgewählte Verarbeitungsschritte, Übergaben und Zustandswirkungen aus Laufartefakten rekonstruierbar; Nachrichten- und Prozessgrenzentests ergänzen diese Belege [D]+[VA] & Keine Vollprüfung aller Informations- und Protokollierungspfade; sichtbarer Kontext belegt keine Nutzung durch das Modell \\
    \shortstack[l]{P1\\(A3)} & Hauptläufe: Full Coverage 0{,}839 / 0{,}495 (Treue/Stress); Korrektheit jeweils 1{,}0, Stützung 1{,}0 / 0{,}957. Drei Stressfall-Ausgaben: Full Coverage 0{,}495--0{,}523 [B] & Nur 43/109 Aussagen in allen drei Ausgaben full; bekannte Quelle, keine allgemeine Stabilitätsaussage. Im Hauptlauf entstanden zwei der 55 Lücken erst nach der Extraktion; teilweise unvollständiger Referenzkontext \\
    \shortstack[l]{P2\\(A3)} & 37 der 55 Stressfall-Lücken mit zuordenbarem Hinweis (Silent-Miss 0{,}327); Treue-Fall 5/5 still [B] & Prüfpotenzial, keine gemessene Behebung; Hinweise verrauscht (18/48 relevant, 1 explizite Warnung); stille Klasse: Detailverlust in verwendeten Units \\
    \shortstack[l]{P3\\(A4)} & 107/109 Klassen unverändert; 2 Verluste aus einer Zusammenführung; 6/7 Merges ohne Klassenverlust [B] & Übergangs-Check prüft Referenz, nicht Inhalt (belegte Kontrolllücke); Verifikation: 12 Paare, nicht verblindet \\
    \shortstack[l]{P4\\(A4/A6)} & Freigabe-Governance am realen Apply-Pfad getestet (4/4); Challenge-Bilanz 6 getestet / 4 teilweise / 2 offen [VA] & Initialer Seed ohne Einzelautorisierung; G01 (nur Warnung) und G04 (Naht-Umgehung) ungetestet; Prüftiefe überwiegend In-Memory (28/36) \\
    \shortstack[l]{P5\\(A5/A7)} & Fünf Falltypen und acht zusätzliche Inhaltsprüfungen: autorisierte Fortschreibung, passende PBI-/Feature-Verbindungen und Ablehnung einer unvollständigen Vorlage fallbezogen belegt [D] & Historische, nicht unabhängige Fälle; teilweise unvollständige Einordnung und Inhaltserhaltung, begrenzte Folgeangleichung; keine Qualitätsquote \\
    \shortstack[l]{P6\\(A3/A5)} & v4: 104/104 Ledger-Ursprungskanten auflösbar; 234/237 Items mit direktem Bezug oder Beitragspfad; 64 aktuelle Einzelreferenzen bei 33 Items offen [VA] & Teilbeitrag ist keine vollständige Herkunft; getrennte Bezugsgrößen, keine semantische Vollprüfung; ein Snapshot \\
    \shortstack[l]{P7\\(A8)} & Drei Issues real erzeugt; ausgewählter Anforderungsinhalt bis zur Projektion feldverglichen; externer Kommentar-Rückweg bis zum Issue-Update; Wiedererkennung abgelehnter Eingänge [D] & kein allgemeiner Drift-Monitor (G08 teilweise); ein Vermerk-Write zunächst fehlgeschlagen \\
    \shortstack[l]{P8\\(A9/A10)} & Wiederaufnahme ohne erneute Ausführung der abgeschlossenen Ledger-Stufe [D]; separater Prozessgrenzen-Smoke mit 14 Prüfungen [VA]; 2 Steward-Aufgaben über reguläre Gates [D] & keine Fortsetzung \emph{innerhalb} laufender Verarbeitung gezeigt; kein Endzustandsvergleich; keine Usability-Aussage \\
    \hline
  \end{tabular}
  \caption[Ergebnisprofil entlang der Prüffragen]{Ergebnisprofil entlang der Prüflandkarte aus
    Tabelle~\ref{t:eval-pruefkarte} --- Befund, Evidenzklasse und
    verbleibende Grenze je Prüffrage mit Bezug zu A1--A10; A1/A2
    querschnittlich ergänzt. Keine pauschale Erfüllungswertung, kein
    Gesamtwert und keine Bestehensschwelle.}
  \label{t:eval-ergebnisprofil}
\end{table}

% Label: t:eval-ergebnisprofil
% Stil: chap6.1-Konvention. P-Nummern = t:eval-pruefkarte (7.1).
% HERKUNFT: Befunde = tex-v2 §§7.3–7.6 (Zahlen: Gold v2/Nenner 109;
% Kanonik runs/w2-goldv2-eval-*.json, z1-stufenanalyse.md rev2,
% fehlerklassen-matrix.md v2, z2-fallblaetter.md rev2,
% traceability-audit.json, Lauf dc711f).
% ============================================================

\begin{table}[htbp]
  \centering
  \small
  \begin{tabular}{>{\raggedright\arraybackslash}p{0.08\textwidth}>{\raggedright\arraybackslash}p{0.40\textwidth}>{\raggedright\arraybackslash}p{0.42\textwidth}}
    \hline
    \textbf{Bezug} & \textbf{Befund (Evidenzklasse)} & \textbf{Verbleibende Grenze} \\
    \hline
    A1/A2 & Ausgewählte Verarbeitungsschritte, Übergaben und Zustandswirkungen aus Laufartefakten rekonstruierbar; Nachrichten- und Prozessgrenzentests ergänzen diese Belege [D]+[VA] & Keine Vollprüfung aller Informations- und Protokollierungspfade; sichtbarer Kontext belegt keine Nutzung durch das Modell \\
    \shortstack[l]{P1\\(A3)} & Hauptläufe: Full Coverage 0{,}839 / 0{,}495 (Treue/Stress); Korrektheit jeweils 1{,}0, Stützung 1{,}0 / 0{,}957. Drei Stressfall-Ausgaben: Full Coverage 0{,}495--0{,}523 [B] & Nur 43/109 Aussagen in allen drei Ausgaben full; bekannte Quelle, keine allgemeine Stabilitätsaussage. Im Hauptlauf entstanden zwei der 55 Lücken erst nach der Extraktion; teilweise unvollständiger Referenzkontext \\
    \shortstack[l]{P2\\(A3)} & 37 der 55 Stressfall-Lücken mit zuordenbarem Hinweis (Silent-Miss 0{,}327); Treue-Fall 5/5 still [B] & Prüfpotenzial, keine gemessene Behebung; Hinweise verrauscht (18/48 relevant, 1 explizite Warnung); stille Klasse: Detailverlust in verwendeten Units \\
    \shortstack[l]{P3\\(A4)} & 107/109 Klassen unverändert; 2 Verluste aus einer Zusammenführung; 6/7 Merges ohne Klassenverlust [B] & Übergangs-Check prüft Referenz, nicht Inhalt (belegte Kontrolllücke); Verifikation: 12 Paare, nicht verblindet \\
    \shortstack[l]{P4\\(A4/A6)} & Freigabe-Governance am realen Apply-Pfad getestet (4/4); Challenge-Bilanz 6 getestet / 4 teilweise / 2 offen [VA] & Initialer Seed ohne Einzelautorisierung; G01 (nur Warnung) und G04 (Naht-Umgehung) ungetestet; Prüftiefe überwiegend In-Memory (28/36) \\
    \shortstack[l]{P5\\(A5/A7)} & Fünf Falltypen und acht zusätzliche Inhaltsprüfungen: autorisierte Fortschreibung, passende PBI-/Feature-Verbindungen und Ablehnung einer unvollständigen Vorlage fallbezogen belegt [D] & Historische, nicht unabhängige Fälle; teilweise unvollständige Einordnung und Inhaltserhaltung, begrenzte Folgeangleichung; keine Qualitätsquote \\
    \shortstack[l]{P6\\(A3/A5)} & v4: 104/104 Ledger-Ursprungskanten auflösbar; 234/237 Items mit direktem Bezug oder Beitragspfad; 64 aktuelle Einzelreferenzen bei 33 Items offen [VA] & Teilbeitrag ist keine vollständige Herkunft; getrennte Bezugsgrößen, keine semantische Vollprüfung; ein Snapshot \\
    \shortstack[l]{P7\\(A8)} & Drei Issues real erzeugt; ausgewählter Anforderungsinhalt bis zur Projektion feldverglichen; externer Kommentar-Rückweg bis zum Issue-Update; Wiedererkennung abgelehnter Eingänge [D] & kein allgemeiner Drift-Monitor (G08 teilweise); ein Vermerk-Write zunächst fehlgeschlagen \\
    \shortstack[l]{P8\\(A9/A10)} & Wiederaufnahme ohne erneute Ausführung der abgeschlossenen Ledger-Stufe [D]; separater Prozessgrenzen-Smoke mit 14 Prüfungen [VA]; 2 Steward-Aufgaben über reguläre Gates [D] & keine Fortsetzung \emph{innerhalb} laufender Verarbeitung gezeigt; kein Endzustandsvergleich; keine Usability-Aussage \\
    \hline
  \end{tabular}
  \caption[Ergebnisprofil entlang der Prüffragen]{Ergebnisprofil entlang der Prüflandkarte aus
    Tabelle~\ref{t:eval-pruefkarte} --- Befund, Evidenzklasse und
    verbleibende Grenze je Prüffrage mit Bezug zu A1--A10; A1/A2
    querschnittlich ergänzt. Keine pauschale Erfüllungswertung, kein
    Gesamtwert und keine Bestehensschwelle.}
  \label{t:eval-ergebnisprofil}
\end{table}


% P10-6: Ergebnisbilanz nicht mit einer einzelnen Zeile nach der Tabelle beginnen.
\begin{samepage}
Das Ergebnisprofil verbindet drei Befundgruppen: Die Ledger-Ausgaben
enthalten überwiegend quellentreue und durch ihre Referenzen gestützte
Claims, erreichen jedoch keine vollständige Abdeckung. Die Zusatzläufe
zeigen trotz ähnlicher Gesamtquoten Unterschiede in den erhaltenen
Inhalten. Zwischenstände und Hinweise machen ausgewählte
Verluste untersuchbar; eine erfolgreiche menschliche Behebung ist
damit nicht nachgewiesen. Kontrolltests und Fallprüfungen belegen
bestimmte Mechanismen der autorisierten Fortschreibung. Die ergänzende
Inhaltsprüfung verfolgt eine neue Entwurfsregel bis zum fachlich
erweiterten Arbeitspaket und eine unvollständige Vorlage bis zu ihrer
begründeten Ablehnung. Daneben bestehen Teilüberlappungen bei der
Einordnung, nicht vollständig erhaltene Altinhalte sowie die bereits
berichteten Grenzen bei Herkunftsanzeige und Folgeangleichung.
\par
\end{samepage}

Die freie Vergleichsbedingung erreicht im Stressfall eine höhere
direkte Abdeckung bei geringerem Tokenaufwand. Sie meldet keine
Lücken im vorgesehenen Kanal. Für die bestehenden Folgestufen fehlen
ihr Dispositionen, Facetten und konventionskonforme Identitäten;
Unit-Referenzen und lauflokale IDs liefert sie. Diese Unterschiede
bleiben getrennt zu bewerten.

Die Bedeutung dieses Ergebnisprofils für den Gestaltungsansatz und
seine Einordnung gegenüber verwandten Arbeiten werden in
Abschnitt~\ref{sec:diskussion-beitrag} diskutiert.

\subsection{Aufwandsprofil}
\label{sec:eval-aufwand}

Tabelle~\ref{t:eval-aufwand} berichtet den Tokenaufwand derselben
vier Hauptläufe und desselben Verarbeitungsausschnitts wie deren
Qualitätsauswertung; für die Zusatzläufe wurde kein neues Aufwandsprofil
erhoben.
Beim Ledger werden alle ausgeführten Stufen samt internen
Wiederholungen summiert. Im Stresslauf blieb der Kanonisierungs-Check
ohne Reparatur-Iteration; ausgeführte Facetten-Reparaturen und der
Referenz-Reparatur-Durchgang sind enthalten. Erhebungsregel und
Stufenzuordnung stehen in Anhang~\ref{anh:evaluation},
\hyperref[anh:hauptlaeufe]{Abschnitt~D.1} und
\hyperref[anh:pruefartefakte]{Abschnitt~D.4}.

Tokens sind technische Aufwandsindikatoren. Kosten, Laufzeit und
menschliche Prüfzeit wurden nicht gemessen; die externe
LLM-Adjudikation zählt nicht zum Systemaufwand. Die vier Hauptläufe
enthielten keine Fehlversuche. Vier Artefakt-Ausfälle der ergänzenden
Modellreihe werden im Anhang gesondert ausgewiesen.

% ============================================================
% Tabelle 7.7-2: Tokenaufwand der bewerteten Läufe — M02 (10.09.2026)
% In Overleaf ablegen als: tables/eval-aufwand.tex
% Aufruf im Fließtext: \input{\tabledir/eval-aufwand.tex}
% Label: t:eval-aufwand
% Stil: chap6.1-Konvention.
% HERKUNFT: thesis-evidence/w2/token-usage-correction-20260910.json.
% B-14: 135.592 Input-Tokens für Ledger-Stresslauf; 46.592 gecachte
% Tokens sind darin enthalten. Übrige sieben Tabellenwerte bestätigt.
% Gleiche Messgrenze wie die Qualitätszahlen (§7.3); Modellreihe inkl.
% 4 F-Ausfällen im Anhang. Erhebungsregel und Korrekturbeleg: D.1.
% ============================================================

\begin{table}[htbp]
  \centering
  \small
  \begin{tabular}{lrrrr}
    \hline
    & \multicolumn{2}{c}{\textbf{Treue-Fall}} & \multicolumn{2}{c}{\textbf{Stressfall}} \\
    & Input & Output & Input & Output \\
    \hline
    Ledger-Kette (alle Stufen) & 26\,731 & 13\,524 & 135\,592 & 41\,133 \\
    Einzel-Agent F & 23\,040 & 5\,708 & 105\,166 & 19\,555 \\
    \hline
  \end{tabular}
  \caption[Tokenaufwand der vier bewerteten Hauptläufe]{Tokenaufwand (Input/Output) der vier bewerteten
    Hauptläufe --- derselbe Verarbeitungsausschnitt wie die
    Qualitätszahlen, beim Ledger einschließlich aller ausgeführten
    Stufen (im Stresslauf: Kanonisierungs-Check ohne
    Reparatur-Iteration; Facetten- und Referenz-Reparaturen
    ausgeführt und enthalten, Anhang~\ref{anh:evaluation}, \hyperref[anh:pruefartefakte]{Abschnitt~D.4}). Technischer
    Indikator, keine Kosten- oder Zeitmessung; Ausfälle der
    ergänzenden Modellreihe gesondert im Anhang.}
  \label{t:eval-aufwand}
\end{table}


Die mehrstufige Kette benötigt im Stressfall rund das 1{,}3-fache
der Eingabe- und das 2{,}1-fache der Ausgabe-Tokens der
Einzel-Agent-Konfiguration. Der höhere Tokenaufwand beschreibt
den Unterschied der untersuchten Gesamtkonfigurationen; welchen
Anteil einzelne Kontroll-, Bilanz- oder Übergabeschritte daran
haben, wurde nicht isoliert bestimmt. Die Governance-Ergebnisse
aus Abschnitt~\ref{sec:eval-kontrollen} stammen zudem aus einer anderen Untersuchung
als diese Tokenmessung und sind nicht deren unmittelbar
gemessener Gegenwert. Eine Verdichtung von Qualität und Aufwand
zu einer Effizienzzahl unterbleibt bewusst.
